% SOURCE_AUTHORITY: sga5_fr_workpass.tex, lines 7841--7873 (LF-normalized unit).
% SOURCE_SHA256: 791F4EFFC5E02832D5D77ED03518C8156D6F07E4C8238B03545DB93D883FBB28
\begin{statement}{Corolário 5.2.5}
Seja
\[
0\longrightarrow X\longrightarrow Y\longrightarrow Z\longrightarrow0
\]
uma sequência exata de \(\Hom(\mathbb N^\circ,\mathcal C)\), com \(J^{n+1}Y_n=0\) para todo \(n\). Supõe-se que:
\begin{enumerate}[label=\textup{(\roman*)}]
\item \(X\) é AR-\(J\)-ádico artiniano (resp. artiniano e noetheriano);
\item \(Z\) é AR-\(J\)-ádico artiniano (resp. noetheriano).
\end{enumerate}
Então \(Y\) é AR-\(J\)-ádico artiniano (resp. noetheriano).
\end{statement}

\subsubsection*{5.3. Sistemas projetivos AR-\(J\)-ádicos noetherianos e \(\delta\)-funtores}

Sejam \(\mathcal C\) e \(\mathcal D\) duas \(A\)-categorias abelianas, e \(F:\mathcal C\to\mathcal D\) um funtor aditivo. O funtor \(F\) se prolonga de maneira evidente a um funtor, denotado também por \(F\),
\[
F:\Hom(\mathbb N^\circ,\mathcal C)\longrightarrow\Hom(\mathbb N^\circ,\mathcal D),
\]
que é aditivo e «comuta com as translações».

\begin{statement}{Proposição 5.3.1}
Sejam \(\mathcal C\) e \(\mathcal D\) duas \(A\)-categorias abelianas e
\[
T=(T^i)_{i\in\mathbb N}
\]
um \(\delta\)-funtor exato \(A\)-linear de \(\mathcal C\) em \(\mathcal D\). Supõe-se que:
\begin{enumerate}[label=\textup{(\roman*)}]
\item Para todo objeto noetheriano \(M\) de \(\mathcal C\), anulado por uma potência de \(J\), e todo inteiro \(n\), \(T^n(M)\) é noetheriano.
\item Para toda \(A\)-álgebra graduada de tipo finito \(B\), anulada por \(J\), e todo \((B,\mathcal C)\)-módulo graduado noetheriano \(N\), os \((B,\mathcal D)\)-módulos graduados \(T^n(N)\) \((n\in\mathbb N)\) são noetherianos.
\end{enumerate}
Então, para todo objeto AR-\(J\)-ádico noetheriano \(X\) de \(\Hom(\mathbb N^\circ,\mathcal C)\), os \(T^n(X)\) são AR-\(J\)-ádicos noetherianos.
\end{statement}
